
\documentclass{article}
\usepackage{colortbl}
\usepackage{makecell}
\usepackage{multirow}
\usepackage{supertabular}

\begin{document}

\newcounter{utterance}

\centering \large Interaction Transcript for game `referencegame', experiment `random\_grids', episode 14 with reflexion{-}Qwen3.5{-}9B.
\vspace{24pt}

{ \footnotesize  \setcounter{utterance}{1}
\setlength{\tabcolsep}{0pt}
\begin{supertabular}{c@{$\;$}|p{.15\linewidth}@{}p{.15\linewidth}p{.15\linewidth}p{.15\linewidth}p{.15\linewidth}p{.15\linewidth}}
    \# & $\;$A & \multicolumn{4}{c}{Game Master} & $\;\:$B\\
    \hline

    \theutterance \stepcounter{utterance}  
    & & \multicolumn{4}{p{0.6\linewidth}}{
        \cellcolor[rgb]{0.9,0.9,0.9}{
            \makecell[{{p{\linewidth}}}]{
                \texttt{\tiny{[P1$\langle$GM]}}
                \texttt{You have attempted this same task before and did not succeed. These are the plans you wrote after those attempts.} \\
\texttt{Trial \#1: New plan: My previous description failed because it contained factual errors regarding the target grid's pattern (the first row has two filled cells separated by a gap, not three consecutive ones) and it was too verbose, potentially missing the most distinctive feature that separates the target from the distractors. The target grid is uniquely defined by having exactly three filled cells in the first row (positions 2, 3, 5) and a specific "staircase" or alternating pattern in the first column (filled, empty, empty, filled, empty). Next time, I will start by identifying the most unique row or column pattern in the target grid that is absent in both distractors, verify the exact coordinates of the filled cells against the visual data, and construct a concise expression focusing solely on these verified unique features without adding unnecessary descriptive clauses.} \\
\\ 
\\ 
\texttt{You are given three grids, where each of them is 5 by 5 in size.} \\
\texttt{Grids have empty cells marked with "▢" and filled cells marked with "X".} \\
\texttt{Your task is to generate a referring expression that best describes the target grid while distinguishing it from the two other distractor grids.} \\
\texttt{The first grid is the target grid, and the following two grids are the distractors.} \\
\\ 
\texttt{Target grid:} \\
\\ 
\texttt{□ X X □ X} \\
\texttt{X □ X X □} \\
\texttt{□ X □ X X} \\
\texttt{X □ X □ X} \\
\texttt{□ X X □ □} \\
\\ 
\texttt{Distractor grid 1:} \\
\\ 
\texttt{□ X □ X □} \\
\texttt{X □ X □ X} \\
\texttt{□ X □ X □} \\
\texttt{X □ X □ X} \\
\texttt{□ X □ X □} \\
\\ 
\texttt{Distractor grid 2:} \\
\\ 
\texttt{□ □ X X □} \\
\texttt{X □ □ X X} \\
\texttt{□ X X □ X} \\
\texttt{X X □ □ X} \\
\texttt{□ X X □ □} \\
\\ 
\texttt{Instruction: Describe the target grid.} \\
\texttt{Generate the referring expression starting with the tag "Expression: " for the given target grid. Omit any other text.} \\
            }
        }
    }
    & \\ \\

    \theutterance \stepcounter{utterance}  
    & \multicolumn{4}{p{0.6\linewidth}}{
        \cellcolor[rgb]{0.9,0.9,0.9}{
            \makecell[{{p{\linewidth}}}]{
                \texttt{\tiny{[P1$\rangle$GM]}}
                \texttt{Expression: The grid has filled cells at (1,2), (1,3), (1,5) in the first row, and a column 2 pattern of filled, empty, filled, empty, filled, distinguishing it from the distractors which have different first{-}row configurations and column patterns.} \\
            }
        }
    }
    & & \\ \\

    \theutterance \stepcounter{utterance}  
    & & & \multicolumn{2}{p{0.3\linewidth}}{
        \cellcolor[rgb]{0.9,0.9,0.9}{
            \makecell[{{p{\linewidth}}}]{
                \texttt{\tiny{[GM$|$GM]}}
                \texttt{Expression: The grid has filled cells at (1,2), (1,3), (1,5) in the first row, and a column 2 pattern of filled, empty, filled, empty, filled, distinguishing it from the distractors which have different first{-}row configurations and column patterns.} \\
            }
        }
    }
    & & \\ \\

    \theutterance \stepcounter{utterance}  
    & & \multicolumn{4}{p{0.6\linewidth}}{
        \cellcolor[rgb]{0.9,0.9,0.9}{
            \makecell[{{p{\linewidth}}}]{
                \texttt{\tiny{[GM$\rangle$P2]}}
                \texttt{You have attempted this same task before and did not succeed. These are the plans you wrote after those attempts.} \\
\texttt{Trial \#1: New plan: I incorrectly identified the "First" grid because I misread the pattern in its second row. The expression requires the second row to start with a filled cell, then a gap, then two consecutive filled cells (pattern: X \_ X X). The "First" grid's second row is "X □ □ X X" (filled, gap, gap, filled, filled), which has two gaps before the consecutive pair. The "Third" grid's second row is "X □ X X □" (filled, gap, filled, filled, gap), which perfectly matches the required pattern. Next time, I will carefully verify the specific sequence of filled cells and gaps in the second row of each candidate grid before making a selection, ensuring the gap is exactly one cell wide between the first filled cell and the pair of consecutive filled cells.} \\
\\ 
\\ 
\texttt{You are given three grids, where each of them is 5 by 5 in size.} \\
\texttt{Grids have empty cells marked with "▢" and filled cells marked with "X".} \\
\texttt{You are also given a referring expression that describes one of the given grids.} \\
\texttt{Your task is to select the grid that matches the given referring expression.} \\
\texttt{Generate only the number (in text) of the grid that matches the given expression by selecting first, second, or third.} \\
\\ 
\texttt{First:} \\
\\ 
\texttt{□ □ X X □} \\
\texttt{X □ □ X X} \\
\texttt{□ X X □ X} \\
\texttt{X X □ □ X} \\
\texttt{□ X X □ □} \\
\\ 
\texttt{Second:} \\
\\ 
\texttt{□ X □ X □} \\
\texttt{X □ X □ X} \\
\texttt{□ X □ X □} \\
\texttt{X □ X □ X} \\
\texttt{□ X □ X □} \\
\\ 
\texttt{Third:} \\
\\ 
\texttt{□ X X □ X} \\
\texttt{X □ X X □} \\
\texttt{□ X □ X X} \\
\texttt{X □ X □ X} \\
\texttt{□ X X □ □} \\
\\ 
\texttt{Expression: The grid has filled cells at (1,2), (1,3), (1,5) in the first row, and a column 2 pattern of filled, empty, filled, empty, filled, distinguishing it from the distractors which have different first{-}row configurations and column patterns.} \\
\texttt{Question: Which grid does the expression refer to?} \\
\texttt{Start with the tag "Answer: ", followed by your selection. Omit any other text.} \\
            }
        }
    }
    & \\ \\

    \theutterance \stepcounter{utterance}  
    & & & \multicolumn{4}{p{0.6\linewidth}}{
        \cellcolor[rgb]{0.9,0.9,0.9}{
            \makecell[{{p{\linewidth}}}]{
                \texttt{\tiny{[GM$\langle$P2]}}
                \texttt{Answer: Second} \\
            }
        }
    }
     \\ \\

    \theutterance \stepcounter{utterance}  
    & & & \multicolumn{2}{p{0.3\linewidth}}{
        \cellcolor[rgb]{0.9,0.9,0.9}{
            \makecell[{{p{\linewidth}}}]{
                \texttt{\tiny{[GM$|$GM]}}
                \texttt{second} \\
            }
        }
    }
    & & \\ \\

    \theutterance \stepcounter{utterance}  
    & \multicolumn{4}{p{0.6\linewidth}}{
        \cellcolor[rgb]{0.95,0.95,0.95}{
            \makecell[{{p{\linewidth}}}]{
                \texttt{\tiny{[P2$\rightarrow$P1]}}
                \texttt{Post{-}episode Reflexion (not sent to GM):} \\
\texttt{New plan: My previous description failed because it hallucinated the pattern of the first row (claiming three consecutive filled cells instead of two separated by one) and misidentified the unique column pattern. The target grid is uniquely defined by having exactly two filled cells in the first row at positions (1,2) and (1,3) with a gap at (1,4), and a specific alternating pattern in the second column (filled, empty, filled, empty, filled). Next time, I will strictly verify the coordinates of the filled cells against the visual data before writing, focusing on the precise count and positions in the first row and the alternating sequence in the second column to ensure the description is factually accurate and distinguishes the target from the distractors.} \\
            }
        }
    }
    & & \\ \\

    \theutterance \stepcounter{utterance}  
    & & & \multicolumn{4}{p{0.6\linewidth}}{
        \cellcolor[rgb]{0.95,0.95,0.95}{
            \makecell[{{p{\linewidth}}}]{
                \texttt{\tiny{[P1$\rightarrow$P2]}}
                \texttt{Post{-}episode Reflexion (not sent to GM):} \\
\texttt{New plan: I incorrectly selected the "Second" grid because I failed to rigorously verify the specific coordinates mentioned in the expression against the actual grid layout. The expression explicitly states filled cells at (1,2), (1,3), and (1,5) in the first row. Upon re{-}examining the grids, the "First" grid has filled cells at (1,3), (1,4), and (1,5), which does not match. The "Second" grid has filled cells at (1,2) and (1,4), missing (1,3). The "Third" grid has filled cells at (1,2), (1,3), and (1,5), which perfectly matches the first{-}row requirement. Additionally, the "Third" grid's column 2 pattern (filled, empty, filled, empty, filled) matches the description, whereas the others do not. Next time, I will systematically cross{-}reference every coordinate listed in the expression with the corresponding grid before making a selection, ensuring no coordinate is overlooked or misaligned.} \\
            }
        }
    }
     \\ \\

\end{supertabular}
}

\end{document}
